\section{Conclusion}\label{sec:conclusion}

In this paper, we aim to reconcile divergent estimates on the returns to rural-urban migration in developing countries by accounting for the role of selection and heterogeneity in returns to migration. Motivated by a multi-period Roy model, we formulate a correlated random coefficient model that allows for location-specific skills and heterogeneous returns.  
We use a person’s migration history as the relevant dimension of heterogeneity, which we interpret as a person’s revealed preference for location and that is conveniently observable in the data. 
Then, we build on \cite{suriSelectionComparativeAdvantage2011}, 
\cite{lemieuxEstimatingEffectsUnions1998}, and   \cite{tjernstromCommentSuri2011} to leverage a linear relationship between returns to migration and comparative advantage that allows us to extrapolate the returns identified from migrant sub-populations to non-migrants.
The group of non-migrants plays a central role in debates on misallocation, whereby non-migrants may reside in areas where they do not live up to their economic potential. 
Moreover, this group is of specific interest to policymakers determining whether to promote migration as a development strategy.

% We take our model to the data which is set in three developing countries, ranging from a low-income country, Tanzania, to a lower middle-income country Indonesia and an upper-middle income country, China. 
We test our model using detailed survey data from three developing countries, Indonesia, China and Tanzania. In line with the existing literature, we confirm the existence of large cross-sectional consumption gaps between rural and urban areas. The gap narrows when we use the panel component of the data and include demographic controls, a time trend, and especially when we include individual fixed effects. 
We then estimate our Group Random Coefficient model and test its assumptions. 
Results show a distinct pattern of the relationship between absolute and comparative advantage that is remarkably consistent across the three countries: individuals with the lowest consumption in rural areas stand most to gain from migrating to urban areas. 
As such, migration can be seen as a pro-poor technology but individuals with the lowest consumption may face barriers to migrate, which might be due to borrowing, liquidity, or information constraints. 
We conclude that individuals are inefficiently sorted across space and that promoting migration would reduce misallocation and, thereby, increase overall growth. 


